\chapter{Marco Teórico y Estado del Arte}
\label{cap:marco}

Este capítulo sitúa el método en el contexto de la predicción de llegadas de transporte público y presenta las
herramientas estadísticas en que se apoya. Se centra en los trabajos y las técnicas que explican de dónde vienen las
ideas del método y por qué se eligieron unas y no otras.

\section{Predicción de la Hora de Llegada}

Predecir cuándo llegará un vehículo a un punto de su recorrido es un problema clásico de los sistemas inteligentes de
transporte. Los enfoques de la literatura pueden agruparse en cuatro familias.

\subsection{Medias Históricas} Se estima el tiempo de recorrido entre dos puntos con el valor típico observado en el
pasado para esa hora y ese día. Son sencillas, estables y fáciles de explicar, pero no reaccionan a lo que le ocurre
al vehículo concreto.

\subsection{Regresión y Modelos de Series Temporales} Relacionan el tiempo de recorrido con variables explicativas
(hora, distancia, número de paradas, demanda) o con su propia historia reciente. Exigen elegir y mantener esas
variables.

\subsection{Filtros de Kalman} Combinan una predicción \emph{a priori} con la observación más reciente, ponderando cada una
por su incertidumbre \citep{kalman1960}. Se han usado para predecir llegadas y salidas a partir de datos de
localización automática de vehículos (AVL) y de conteo automático de viajeros (APC) \citep{shalaby2004}.

\subsection{Redes Neuronales} Aprenden la relación entre la posición y el tiempo restante a partir de muchos datos
\citep{chien2002,jeong2005}. Suelen ser precisas cuando hay abundancia de datos homogéneos, pero son difíciles de
auditar y de explicar a quien usa el sistema.

\subsection{Trabajos Recientes} Dos revisiones recientes recogen el estado del campo: \citet{kumar2025review}, que cubre
el proceso completo (datos, modelos, discretización del recorrido y medidas de evaluación), y
\citet{singh2022review}, centrada en los modelos de inteligencia artificial. Partir el viaje en tramos entre paradas
consecutivas y predecir cada uno por separado es habitual \citep{gal2017journey}. La variabilidad del tiempo de
recorrido depende mucho de la hora del día \citep{mazloumi2010variability}, lo que justifica aprender por franjas.
\citet{lin1999experimental} compararon algoritmos que combinan el GPS con el horario en un servicio rural de baja
frecuencia, un caso parecido al de los autobuses interurbanos. Entre los filtros de Kalman, \citet{cathey2003prescription} propusieron el esquema de
seguimiento, filtrado y predicción para datos AVL, y \citet{vanajakshi2009kalman} lo aplicaron a autobuses en tráfico
heterogéneo. Entre las redes neuronales se han usado redes recurrentes \citep{pang2019rnn}, redes LSTM convolucionales
que predicen a la vez todos los enlaces de una línea \citep{petersen2019convlstm}, redes de grafos con atención para
datos escasos \citep{ma2022multiattention} y transformadores entrenados con datos abiertos
\citep{jeong2024battransformer}. En la predicción del tráfico, el aprendizaje profundo no siempre es la mejor técnica,
y los métodos de referencia sencillos y bien ajustados siguen siendo importantes \citep{manibardo2022deep}.

Otros trabajos parten, como este, de los datos que publican los operadores. \citet{wessel2017retrospective}
reconstruyeron el horario realmente cumplido a partir de un feed público de posiciones, igual que aquí se reconstruye
la verdad de referencia. \citet{cats2016realtime} evaluaron las predicciones que ya se mostraban a los viajeros en
Estocolmo y encontraron que subestimaban la espera. \citet{garciamaurino2024improving} corrigieron los errores
sistemáticos de las predicciones de los operadores de Madrid y París solo con datos de sus API públicas.
\citet{woodard2017predicting} modelaron, para viajes en coche, un efecto propio de cada viaje y recurrieron a
parámetros agrupados cuando un enlace tenía pocas observaciones, dos recursos análogos a la corrección por viaje y a la
contracción de este método.

El método de esta memoria pertenece a la primera familia en su forma de estimar tiempos (valores típicos por tramo y
franja) y combina esa estimación con la observación en tiempo real de la posición, que hace de «reinicio» de la predicción en cada
actualización del feed. La elección es deliberada: con un dato de posición por minuto y sin retrasos publicados, un
modelo explicable, robusto y medible tramo a tramo resultó más útil que uno de caja negra. Las ideas que se
probaron y no mejoraron la medida (entre ellas, corregir con la velocidad del último minuto del vehículo, la forma más
simple de filtrar) se recogen en la sección «\nameref{sec:descartadas}». Un filtro de Kalman sobre la desviación del
propio autobús en los tramos ya terminados, suavizada y olvidada con el tiempo, sí reduce un 4\,\% el error medio de las
sumas de tramos (sección «\nameref{sec:comparacion}»): lo que decide es qué se filtra y cuánto se recuerda
(sección «\nameref{sec:otros-enfoques}»).

\section{Datos de Transporte: GTFS y SIRI}

GTFS \citep{gtfs} es el formato \emph{de facto} para publicar horarios programados. Describe agencias, líneas
(\texttt{routes}), viajes (\texttt{trips}), paradas (\texttt{stops}), horas de paso (\texttt{stop\_times}),
calendarios (\texttt{calendar}, \texttt{calendar\_dates}) y trazados (\texttt{shapes}). Dos campos son centrales
aquí: \texttt{shape\_dist\_traveled}, la distancia recorrida a lo largo del trazado hasta cada parada y cada punto del
trazado, y \texttt{timepoint}, que indica si la hora de paso es exacta (\texttt{1}) o aproximada por interpolación
(\texttt{0}).

GTFS Realtime \citep{gtfsrt} y SIRI \citep{siri} son los formatos más usados para datos en tiempo real. SIRI Vehicle
Monitoring (SIRI-VM) publica, para cada vehículo, la actividad de su viaje: referencia del viaje
(\texttt{VehicleJourneyRef}), posición (\texttt{VehicleLocation}), hora de la posición
(\texttt{LocationRecordedAtTime}) y llamada a la próxima parada (\texttt{MonitoredCall}), con su orden
(\texttt{Order} o \texttt{VisitNumber}).

\section{Ajuste al Mapa}

Situar una posición GPS sobre una red o un recorrido se conoce como \emph{map matching} \citep{quddus2007}. Los
métodos van desde la proyección geométrica sobre el segmento más próximo hasta modelos probabilísticos que tienen en
cuenta la topología y la trayectoria. Aquí el recorrido se conoce de antemano (es el trazado del viaje
que el propio feed declara), lo que reduce el problema a una proyección sobre una polilínea. La dificultad que queda
es la ambigüedad cuando el recorrido pasa dos veces por la misma calle; el capítulo «\nameref{cap:localizacion}» la resuelve
con una ventana alrededor de la próxima parada y una regla de desempate.

Para distancias de unos pocos kilómetros basta una proyección equirrectangular local \citep{snyder1987}, que trata la
superficie terrestre como un plano alrededor del punto de interés.

\section{Estadística Robusta}

La media de una muestra es muy sensible a valores extremos: un único autobús parado veinte minutos por un accidente
desplaza la media de un tramo tanto como muchos autobuses normales. Los estimadores \emph{robustos}
\citep{huber1964,tukey1977} limitan esa influencia. La mediana es el más sencillo: su punto de ruptura es del 50\,\%,
es decir, la mitad de la muestra tendría que estar contaminada para llevarla a un valor arbitrario. El método usa
medianas en todos sus agregados, y además toma un solo valor por vehículo y viaje, para que un vehículo lento visto
muchas veces no pese más que uno visto una vez.

\section{Contracción y Modelos Jerárquicos}
\label{sec:marco-contraccion}

Cuando se estiman muchas medias a la vez (un tiempo para cada tramo y cada franja horaria), algunas se apoyan en muy
pocas observaciones y son ruidosas. \citet{jamesstein1961} demostraron que, al estimar simultáneamente tres o más medias de
distribuciones normales, un estimador que contrae las medias observadas hacia un punto fijo tiene un error cuadrático
total esperado menor que el de usar cada media por separado. \citet{efronmorris1975} mostraron su utilidad con datos
reales, contrayendo cada media hacia la media común del grupo. En términos bayesianos
\citep{gelman2013}, la contracción equivale a un modelo jerárquico con \emph{agrupamiento parcial}: cada grupo tiene su
propio parámetro, pero estos se parecen entre sí, y la estimación de cada grupo es una media ponderada de su dato
propio y del valor común, en la que el peso del dato propio crece con el número de observaciones del grupo.

El método aplica esta idea en dos sitios: el tiempo de un tramo en una franja se contrae hacia el tiempo del tramo a
cualquier hora (sección «\nameref{sec:contraccion-tramo}»), y el factor de ritmo se contrae hacia~1 (sección «\nameref{sec:k}»).

\section{Validación}
\label{sec:marco-validacion}

Medir el error con los mismos datos con que se aprendió sobreestima la precisión. La \emph{validación cruzada}
\citep{kohavi1995,hastie2009} separa los datos en particiones, aprende con unas y mide con otras. Cuando los datos
tienen orden temporal, lo honesto es aprender con el pasado y medir con el futuro (\emph{validación con origen
móvil}, \citealp{hyndman2021}). El método usa ambas: una partición por recorridos (un autobús en un viaje de un día) para
decidir si se escribe una nueva calibración y una partición por días para comparar reglas. La primera comparte los
días entre aprendizaje y prueba, y con ellos el tráfico y la meteorología de cada día, por lo que las reglas nuevas se
deciden con la segunda. \citet{hyndmankoehler2006} distinguen entre las medidas de
error que dependen de la escala de los datos, como el error absoluto medio, y las que no: las primeras sirven para
comparar métodos sobre unos mismos datos, pero no entre conjuntos de datos con escalas distintas. Aquí todas las
predicciones se miden en la misma unidad (segundos), de modo que se usan errores absolutos y porcentajes de aciertos
dentro de umbrales que tienen sentido para quien espera el autobús.
